% Part: incompleteness
% Chapter: introduction
% Section: historical-background

\documentclass[../../../include/open-logic-section]{subfiles}

\begin{document}

% SEGMENT OLP-0276-S01

\olfileid{inc}{int}{bgr}

\olsection{வரலாற்றுப் பின்னணி}

முழுமையின்மைத் தேற்றங்களை அவற்றின் சூழலில் புரிந்துகொள்ள உதவும் வரலாற்று
வளர்ச்சிகளை இப்பிரிவில் சுருக்கமாக விவாதிப்போம். குறிப்பாக, கணிதத்
தருக்கத்தின் வரலாற்றை மிகக் கோட்டளவாக மேலோட்டம் செய்தபின், கணிதத்தின்
அடித்தளங்களின் வரலாறு பற்றிச் சில சொற்கள் கூறுவோம்.

\begin{digress}
  ``கணிதத் தருக்கம்'' என்ற சொற்றொடர் இருபொருள்படக்கூடியது. ``கணிதம்'' என்ற
சொல் பொருளடக்கத்தை விவரிப்பதாகக் கொண்டு, கணிதச் சிந்தனையின் கொள்கைகளைக்
குறிக்கும் ``கணிதத்தின் தருக்கம்'' என்று பொருள்கொள்ளலாம். அல்லது அது
முறைகளை விவரிப்பதாகக் கொண்டு, சிந்தனைக் கொள்கைகளை கணிதமுறையில் ஆராய்வதைக்
குறிக்கும் ``தருக்கத்தின் கணிதம்'' என்று பொருள்கொள்ளலாம். பின்வரும் வரலாறு
கணிதத் தருக்கத்தை இவ்விரு பொருள்களிலும், பெரும்பாலும் ஒரே நேரத்தில்,
உள்ளடக்குகிறது.
\end{digress}

ஏறத்தாழ 384--322~\textsc{கி.மு.} காலத்தில் வாழ்ந்த அரிஸ்டாட்டிலிடமே
தருக்கத்தின் ஆய்வு அடிப்படையில் தொடங்கியது. அவரது \emph{வகைப்பாடுகள்},
\emph{முந்தைய பகுப்பாய்வுகள்}, \emph{பிந்தைய பகுப்பாய்வுகள்} ஆகிய நூல்கள்,
முக்கூற்று வாதத்தைப் பற்றிய ஆழமான முறைப்படியான ஆய்வு உட்பட, அறிவியல்
சிந்தனையின் கொள்கைகளை முறைப்படி ஆராய்கின்றன.

இடைக்காலம் முழுவதும் பள்ளிவழித் தத்துவத்தில் அரிஸ்டாட்டிலின் தருக்கமே
மேலோங்கியது. பதினெட்டாம் நூற்றாண்டின் பிற்பகுதியில்கூட, அரிஸ்டாட்டிலின்
தருக்கம் முழுநிறைவானது என்றும் அதற்குத் திருத்தம் தேவையில்லை என்றும் காந்த்
வாதிட்டார். ஆனால் கணிதச் சிந்தனையின் மிக மேலோட்டமான கூறுகளைத் தவிர வேறெதையும்
மாதிரிப்படுத்த முடியாத அளவுக்கு முக்கூற்று வாதக் கோட்பாடு வரம்புடையது. அதற்கு
ஒரு நூற்றாண்டுக்கு முன்பு, நியூட்டனின் சமகாலத்தவரான லைப்னிட்ஸ், தருக்கச்
சிந்தனைக்கான ஒரு முழுமையான ``கணிப்புமுறையை'' கற்பனை செய்தார். அத்தகைய
கணிப்புமுறையை வடிவமைப்பதற்கான சில தொடக்கப் படிகளை எடுத்து, கூற்றுத்
தருக்கத்தின் ஒரு வடிவத்தைச் சாரமாக விவரித்தார்.

பத்தொன்பதாம் நூற்றாண்டு தருக்கத்திற்கு ஒரு திருப்புமுனையாக அமைந்தது. 1854
இல் ஜார்ஜ் பூல் எழுதிய \emph{சிந்தனையின் விதிகள்}, நவீன விளக்கங்களிலிருந்து
அதிகம் வேறுபடாத கூற்றுத் தருக்கத்தின் விரிவான இயற்கணித ஆய்வைக் கொண்டது.
1879 இல் கோட்லோப் ஃபிரேகே தமது \emph{பெக்ரிஃப்ஸ்ஷ்ரிஃப்ட்}
(கருத்தெழுத்து) நூலை வெளியிட்டார். அது கூற்றுத் தருக்கத்தை அளவையடைகளாலும்
தொடர்புகளாலும் விரிவாக்கி, முதல்தர தருக்கத்தை உள்ளடக்குகிறது. உண்மையில்,
ஃபிரேகேயின் தருக்க முறைமைகள் உயர்தரத் தருக்கத்தையும் அதற்கு மேலும் பலவற்றையும்
உள்ளடக்கின. தமது \emph{எண்கணிதத்தின் அடிப்படை விதிகள்} நூலில், முற்றிலும்
தருக்கக் கருதுகோள்களிலிருந்து தமது பெக்ரிஃப்ஸ்ஷ்ரிஃப்ட் முறைமையில் எண்கணிதம்
முழுவதையும் வருவிக்க முடியும் என்று காட்ட ஃபிரேகே முற்பட்டார். ஆனால் 1902
இல் ரசல் காட்டியபடி, இக்கருதுகோள்கள் முரணுடையவையாகத் தெரிந்தன. அந்த
முரணுடைய அடிகோளை ஒதுக்கிவிட்டால், ஃபிரேகே ஏறத்தாழத் தனியொருவராக நவீனத்
தருக்கத்தைக் கண்டுபிடித்தார்; இது வியப்பூட்டும் சாதனை. பூலுக்குப் பிந்தைய,
இயற்கணித நோக்குடைய பியர்ஸ், ஷ்ரோடர் உள்ளிட்ட சிந்தனையாளர்களும் அளவையடைத்
தருக்கத்தைத் தனித்துவமாக உருவாக்கினர்.

இப்போது கணிதத்தின் அடித்தளங்களில் ஏற்பட்ட வளர்ச்சிகளுக்கு வருவோம். கணிதத்தில்
தருக்கம் முக்கியப் பங்காற்றுவதால், இவ்வளர்ச்சிகளுக்கும் இப்போது விவரித்த
வளர்ச்சிகளுக்கும் இடையே இயல்பாகவே மிகுந்த இடைவினை உள்ளது. எடுத்துக்காட்டாக,
தமது தருக்கக் கட்டமைப்பை மட்டும் அடிப்படையாகக் கொண்டு கணிதம் முழுவதையும்
அமைக்க முடியும் என்று காட்டும் வெளிப்படையான நோக்கத்துடனேயே ஃபிரேகே தமது
தருக்கத்தை உருவாக்கினார். குறிப்பாக, காந்த் கூறியதுபோல் கணிதம் அனுபவமுன்
\emph{தொகுப்பாய்வு} மெய்மைகளைக் கொண்டிராமல், அனுபவமுன் \emph{பகுப்பாய்வு}
மெய்மைகளைக் கொண்டுள்ளது என்று காட்ட அவர் விரும்பினார்.

கணிதம் என்ற துறை முறையாகக் கிரேக்கர்களிடமே பிறந்தது என்று பலர் கருதுகின்றனர்.
கி.மு. 300 அளவில் எழுதப்பட்ட யூக்ளிடின் \emph{மூலக்கூறுகள்}, கண்டிப்பையும்
துல்லியத்தையும் வலியுறுத்தும், முதிர்ந்த கிரேக்கக் கணிதத்தின் பிரதிநிதியாக
ஏற்கெனவே விளங்கியது. யூக்ளிடின் \emph{மூலக்கூறுகள்} நூலிலுள்ள வரையறைகளும்
நிறுவல்களும் இன்றைய மேல்நிலைப்பள்ளி வடிவியல் பாடநூல்களில் ஏறத்தாழ மாறாமல்
நிலைத்துள்ளன (மேல்நிலைப்பள்ளிகளில் இன்னும் வடிவியல் கற்பிக்கப்படும் அளவுக்கு).
கணிதத்தில் மட்டுமன்றி தத்துவத்தின் கிளைகளிலும் கண்டிப்பான வாதத்துக்கான
முன்மாதிரியாக இக்கணிதச் சிந்தனை மாதிரி கருதப்பட்டது. (ஸ்பினோசா அறவியல்,
சமய வாதங்களைக்கூட யூக்ளிடிய முறையில் முன்வைத்தார்; அதைக் காண்பது
விந்தையாக உள்ளது!)

பதினேழாம் நூற்றாண்டில் நியூட்டனும் லைப்னிட்ஸும் நுண்கணிதத்தைக்
கண்டுபிடித்தனர். (முதன்மை உரிமை பற்றிய கடுமையான சர்ச்சை நூற்றாண்டுகளாக
நீடித்தது; ஆனால் இரு வளர்ச்சிகளும் பெரும்பாலும் தனித்தனியே ஏற்பட்டன என்று
இன்றைய அறிஞர்கள் பலர் கருதுகின்றனர்.) முடிவிலியாகச் சிறிய அளவுகளின் முடிவிலித்
தொகைகள் போன்றவற்றைப் பற்றிய சிந்தனை நுண்கணிதத்தில் இடம்பெறுகிறது. இக்கூறுகள்,
கடவுள் நம்பிக்கை தமது காலக் கணிதத்தைவிடக் குறைந்த பகுத்தறிவுடையதல்ல என்று
வாதிட்ட பிஷப் பெர்க்லியின் விமர்சனத்தைத் தூண்டின. பதினெட்டாம் நூற்றாண்டில்
நுண்கணித முறைகள் பரவலாகப் பயன்படுத்தப்பட்டன. எடுத்துக்காட்டாக, லியோனார்ட்
ஆய்லர் முடிவிலித் தொகைகள் அடங்கிய கணக்கீடுகளைக் கொண்டு வியப்பூட்டும்
முடிவுகளைப் பெற்றார்.

பத்தொன்பதாம் நூற்றாண்டில் கணிதவியலாளர்கள் நுண்கணிதத்தை உறுதியான அடித்தளத்தில்
அமைத்து பெர்க்லியின் விமர்சனங்களை எதிர்கொள்ள முயன்றனர். கோஷி, வையர்ஸ்ட்ராஸ்,
போல்சானோ முதலியோரின் முயற்சிகள், ``எப்சிலன்கள், டெல்டாக்கள்'' வழியாக
எல்லை, தொடர்ச்சி, வகையீடு, தொகையீடு ஆகியவற்றுக்கான நமது இன்றைய
வரையறைகளுக்கு வழிவகுத்தன; வேறு சொற்களில், நுண்ணளவுகளைச் சுட்டாமல் அவை
வரையறுக்கப்பட்டன. அதே நூற்றாண்டின் பிற்பகுதியில் கணிதவியலாளர்கள் மேலும்
சென்று, மெய்யெண்கள் உட்பட நுண்கணிதத்தின் எல்லாக் கூறுகளையும் இயல் எண்களின்
வழியில் விளக்க முயன்றனர். (குரோனெக்கர்: ``முழு எண்களைக் கடவுள் படைத்தார்;
மற்ற அனைத்தும் மனிதனின் படைப்பு.'') 1872 இல் டெடகிண்ட் எழுதிய
``தொடர்ச்சியும் விகிதமுறா எண்களும்'' என்ற கட்டுரையில், விகிதமுறு எண்களின்
கணங்களாக மெய்யெண்களை எவ்வாறு ``கட்டமைப்பது'' என்று காட்டினார் (விகிதமுறு
எண்களை இயல் எண்களின் ஜோடிகளாகக் கருதலாம் என்பதை நீங்கள் அறிவீர்கள்).
1888 இல் அவர் ``வாஸ் ஸின்ட் உண்ட் வாஸ் ஸொலன் டீ சாலன்'' (ஏறத்தாழ,
``இயல் எண்கள் என்றால் என்ன, அவை எவ்வாறு இருக்க வேண்டும்?'') என்ற நூலை
எழுதி, முற்றிலும் ``தருக்க'' வழியில் இயல் எண்களை விளக்க முற்பட்டார். 1887
இல் குரோனெக்கர் ``ஊபர் டென் சால்பெக்ரிஃப்'' (``எண் என்ற கருத்து பற்றி'')
எழுதி, எல்லாக் கணிதப் பொருள்களையும் முழுக்களின் வழியாகக் குறிப்பிடுவதைப்
பற்றிப் பேசினார். 1889 இல் ஜூசெப்பே பியானோ இயல் எண்களுக்கான முறைப்படியான
குறியீட்டு அடிகோள்களை வழங்கினார்.

பத்தொன்பதாம் நூற்றாண்டின் முடிவு, முடிவிலியை அணுகுவதிலும் புதிய துணிவைக்
கொண்டுவந்தது. அதற்கு முன்பு, முடிவிலிப் பொருள்களும் கட்டமைப்புகளும் (இயல்
எண்களின் கணம் போன்றவை) மிகுந்த கவனத்துடன் கையாளப்பட்டன. ``முடிவிலியாகப்
பல'' என்பது ``நீங்கள் விரும்பும் அளவுக்குப் பல'' என்றும், ``எல்லையில்
நெருங்குகிறது'' என்பது ``நீங்கள் விரும்பும் அளவுக்கு அருகில் வருகிறது''
என்றும் புரிந்துகொள்ளப்பட்டது. ஆனால் முடிவிலியை உள்ளவாறே ஏற்க முடியும்
என்று ஜார்ஜ் காண்டோர் காட்டினார். காண்டோர், டெடகிண்ட் முதலியோரின் பணி,
இப்போது பரவலாக ஏற்கப்படும் கணிதத்தின் பொதுவான கணக்கோட்பாட்டு புரிதலை
அறிமுகப்படுத்த உதவியது.

இது இருபதாம் நூற்றாண்டில் தருக்கத்திலும் அடித்தளங்களிலும் ஏற்பட்ட
வளர்ச்சிகளுக்கு நம்மைக் கொண்டுவருகிறது. 1902 இல் ஃபிரேகேயின் தருக்க
முறைமையிலுள்ள முரண்பாட்டை ரசல் கண்டுபிடித்தார். 1904 இல், ``தேர்வு
அடிகோள்'' எனப்படும் அடிகோளைப் பயன்படுத்தி, காண்டோரின் நன்கு-வரிசைப்படுத்தல்
கொள்கையை செர்மெலோ நிறுவினார்; இவ்வடிகோளின் நியாயத்தன்மை பெரும் விவாதத்தைத்
தூண்டியது. 1910 முதல் 1913 வரை ரசல், ஒயிட்ஹெட் ஆகியோரின்
\emph{பிரின்சிபியா மேத்தமாட்டிக்கா}வின் மூன்று தொகுதிகள் வெளிவந்து,
தருக்க அடித்தளத்தில் கணிதத்தை நிறுவும் ஃபிரேகேயின் திட்டத்தை விரிவாக்கின.
ஆனால் முற்றிலும் தருக்கக் கொள்கைகளாக நியாயப்படுத்தக் கடினமாகத் தோன்றிய இரு
கொள்கைகளை ரசலும் ஒயிட்ஹெட்டும் ஏற்க வேண்டியிருந்தது: முடிவிலி அடிகோளும்
``குறைக்கத்தக்கதன்மை'' அடிகோளும். 1900களில் கணிதத்தில் ``சுயஉள்ளடக்க
வரையறைகள்'' பயன்படுத்தப்படுவதை புவான்காரே விமர்சித்தார். 1910களில்,
விலக்கப்பட்ட நடுவின் விதியை~($!A \lor \lnot !A$) பயன்படுத்தாத ஓர்
``உள்ளுணர்விய'' அடித்தளத்தில் கணிதம் முழுவதையும் மீளமைக்க புரோவர்
முன்மொழியத் தொடங்கினார்.

உண்மையிலேயே விசித்திரமான காலம்!{} கணிதம் முழுவதையும் தருக்கமாகக் குறைக்கும்
திட்டம் இப்போது ``தருக்கவாதம்'' எனப்படுகிறது; மேலே கூறிய இடர்ப்பாடுகளால்
அது தோல்வியடைந்ததாகப் பொதுவாகக் கருதப்படுகிறது. உள்ளுணர்விய மனக்
கட்டமைப்புகளின் வழியாகக் கணிதத்தை உருவாக்கும் திட்டம் ``உள்ளுணர்வியம்''
எனப்படுகிறது; நடைமுறைக் கணிதத்தின் மீது அது அளவுக்கு மீறிய கடுமையான
கட்டுப்பாடுகளை விதிப்பதாகக் கருதப்படுகிறது. நூற்றாண்டுத் திருப்பத்தில்,
எக்காலத்திலும் மிகுந்த செல்வாக்குடைய கணிதவியலாளர்களுள் ஒருவரான டேவிட்
ஹில்பர்ட், காண்டோரும் டெடகிண்டும் அறிமுகப்படுத்திய புதிய அருவமான முறைகளை
வலுவாக ஆதரித்தார்: ``காண்டோர் நமக்காக உருவாக்கிய சொர்க்கத்திலிருந்து எவரும்
நம்மை விரட்ட முடியாது.'' அதே நேரத்தில், இப்புதிய முறைகளின் அடித்தளம் பற்றிய
விமர்சனங்களையும் அவர் உணர்ந்திருந்தார் (விந்தையாக, அவை இப்போது
``செவ்வியல்'' எனப்படுகின்றன). இரு நோக்கங்களையும் ஒருசேர அடையும் வழியை அவர்
முன்மொழிந்தார்:
\begin{enumerate}
\item செவ்வியல் முறைகளை முறைப்படியான அடிகோள்களாலும் விதிகளாலும்
  பிரதிநிதித்துவப்படுத்துக; ஓர் அடிகோள் முறைமையில் கணிதக் கேள்விகளை
  !!{formula}s ஆகப் பிரதிநிதித்துவப்படுத்துக.
\item இம்முறைப்படியான கழித்தறிதல் முறைமைகள் முரணற்றவை என்று நிறுவப்
  பாதுகாப்பான, ``முடிவுறுமுறை'' முறைகளைப் பயன்படுத்துக.
\end{enumerate}

ஹில்பர்டின் பணி முதல் நோக்கத்தை அடைவதில் பெருந்தொலைவு சென்றது. 1899 இல்
தமது புகழ்பெற்ற \emph{வடிவியலின் அடித்தளங்கள்} நூலில் வடிவியலுக்காக இதைச்
செய்திருந்தார். அடுத்தடுத்த ஆண்டுகளில், வடிவியலுக்காக ஹில்பர்ட் செய்ததைப் பிற
கணிதத் துறைகளுக்குச் செய்ய அவரும் அவரது மாணவர்கள், உடன்பணியாளர்கள் பலரும்
உழைத்தனர். எண்கணிதத்திற்கும் பகுப்பாய்விற்கும் ஹில்பர்டே அடிகோள் முறைமைகளை
வழங்கினார். செர்மெலோ கணக் கோட்பாட்டை அடிகோளாக்கினார்; ஃபிராங்கல், ஸ்கோலம்,
ஃபான் நியூமன் முதலியோர் அதை விரிவாக்கினர். 1920களின் நடுப்பகுதிக்குள்,
கணிதம் ``முழுவதற்குமான'' அடிகோளாக்கம் என்று உரிமை கோரிய இரு அணுகுமுறைகள்
இருந்தன: ரசல், ஒயிட்ஹெட் ஆகியோரின் \emph{பிரின்சிபியா மேத்தமாட்டிக்கா},
மற்றும் செர்மெலோ--ஃபிராங்கல் கணக் கோட்பாடு என அறியப்பட்ட முறைமை.

1921 இல், இம்முறைமைகள் முரணற்றவை என்று நிறுவும் நோக்கத்தை அடைவதற்கான ஓர்
ஆராய்ச்சித் திட்டத்தை ஹில்பர்ட் தொடங்கினார். அவரது மாணவர்களில் குறிப்பாக
பெர்னாய்ஸ், ஆக்கர்மன், பின்னர் கென்ட்சன் ஆகியோர் இத்திட்டத்தில் உதவினர்.
இந்நோக்கத்தை அடைவதற்கான அடிப்படை எண்ணம், கணிதத்தில் ஒரு முரண்பாட்டை
!!{derivation} செய்ய முடியுமா என்ற கேள்வியை, குறிகளின் சாத்தியத் தொடர்கள்
பற்றிய சேர்வியல் சிக்கலாக மாற்றுவதாகும். அதாவது எண்கணிதம், பகுப்பாய்வு
அல்லது கணக் கோட்பாட்டிற்கான அடிகோள் முறைமையின் அடிகோள்களிலிருந்து,
எடுத்துக்காட்டாக~$!A \land \lnot !A$ ஐச் சரியாக !!{derivation} செய்யும்
வரையறையை நிறைவுசெய்யும் வாக்கியத் தொடர்கள் இருக்க முடியுமா என்ற சிக்கலாக
மாற்றுவதாகும். அத்தகைய குறித்தொடர் இருக்க முடியாது என்பதற்கான நிறுவல் தானும்
ஒரு கணித நிறுவலாக இருப்பதால், இந்த அடிகோள் முறைமைகளில் அதை முறைப்படுத்த
முடியும். வேறு சொற்களில், எடுத்துக்காட்டாக எண்கணிதம் முரணற்றது என்று கூறும்
$\OCon$ என்ற வாக்கியம் ஒன்று இருக்கும். மேலும், அதன் நிறுவலுக்கு மிகவும்
கட்டுப்பட்ட ``முடிவுறுமுறைச்'' சாதனங்கள் மட்டுமே தேவைப்பட்டால் குறிப்பாக,
குறித்த முறைமைகளில் இவ்வாக்கியம் வருவிக்கத்தக்கதாக இருக்க வேண்டும்.

உருவாக்கப்பட்ட அடிகோள் முறைமைகள் ஒவ்வொரு கணிதக் கேள்வியையும் தீர்க்க வேண்டும்
என்ற இரண்டாம் நோக்கத்தை இரு வழிகளில் துல்லியமாக்கலாம். ஒரு வழியில் அதை
இவ்வாறு கூறலாம்: கணிதத்திற்கான ஓர் அடிகோள் முறைமையின் மொழியிலுள்ள எந்த
வாக்கியம்~$!A$ க்கும், $!A$ அல்லது $\lnot !A$ அடிகோள்களிலிருந்து
வருவிக்கத்தக்கதாக இருக்கும். இது உண்மையானால், அடிகோள்களின் அடிப்படையில்
நிறுவவும் மறுக்கவும் முடியாத வாக்கியங்களும், அடிகோள்கள் தீர்க்காத கேள்விகளும்
இருக்காது. இப்பண்புள்ள அடிகோள் முறைமை \emph{நிறைவானது} எனப்படுகிறது. எந்தக்
குறிப்பிட்ட வாக்கியத்திற்கும் இரு மாற்றுகளில் எது நிலவுகிறது என்று
தீர்மானிப்பது இன்னும் கடினமான பணியாக இருக்கலாம். ஆனால் கொள்கையளவில் அதைச்
செய்ய ஒரு முறை இருக்க வேண்டும். ஹில்பர்ட் கருதிய அடிகோள் மற்றும்
!!{derivation} முறைமைகளில், நிறைவு அத்தகைய முறை இருப்பதை உண்மையில்
உட்படுத்தும்---ஆனால் இதை ஹில்பர்ட் உணரவில்லை. கேள்வியைப் புரிந்துகொள்வதற்கான
இரண்டாம் வழி பின்வரும் வலுவான தேவை: கொடுக்கப்பட்ட வாக்கியம்~$!A$ அடிகோள்களிலிருந்து
வருவிக்கத்தக்கதா இல்லையா என்பதைத் தீர்மானிக்கும் இயந்திரமுறையான, கணக்கீட்டு
முறை ஒன்று இருக்க வேண்டும்.

1931 இல், இத்திட்டம் வெற்றியடைய முடியாது என்று காட்டிய இரு
``முழுமையின்மைத் தேற்றங்களை'' கோடல் நிறுவினார். கணிதத்திற்கான எந்த அடிகோள்
முறைமையும் நிறைவானதல்ல; குறிப்பாக, அடிகோள்களின் முரணின்மையை வெளிப்படுத்தும்
வாக்கியத்தை நிறுவவோ மறுக்கவோ முடியாது.

இது ஹில்பர்டின் மூலத் திட்டத்துக்குப் பேரடியைக் கொடுத்தது. ஆனால் கணிதத்தில்
அடிக்கடி நிகழ்வதுபோல், ஆராய்ச்சிக்கான ஆர்வமூட்டும் புதிய வழிகளையும் அது
திறந்தது. கணிதம் முழுவதையும் உள்ளடக்கும் ஒரே முறைப்படியான முறைமை இல்லையெனில்,
குறுகிய எல்லையுள்ள முறைமைகளை உருவாக்கி அவற்றில் எதை நிறுவலாம் என்று
ஆராய்வது பொருத்தமானது. இம்முறைமைகளின் முரணின்மையை நிறுவக் குறைவாகக்
கட்டுப்பட்ட நிறுவல் முறைகளை உருவாக்குவதும், அவற்றின் முரணின்மையை நிறுவுவது
எவ்வளவு கடினம் என்பதை அளவிடுவதற்கான வழிகளைக் கண்டறிவதும் பொருத்தமானவை.
(ஏறத்தாழ) ஒவ்வொரு முறைப்படியான முறைமையிலும் அதனால் தீர்க்க முடியாத கேள்விகள்
உள்ளன என்று கோடல் காட்டியதால், குறிப்பிட்ட முறைமை தீர்க்க முடியாத
``ஆர்வமூட்டும்'' கேள்விகளைத் தேடுவதும், அவற்றைத் தீர்க்க ஒரு முறைப்படியான
முறைமை எவ்வளவு வலுவாக இருக்க வேண்டும் என்று கண்டறிவதும் பொருத்தமானவை.
இன்றுவரை, தருக்கவியலாளர்கள் இக்கேள்விகளை நிறுவல் கோட்பாடு என்ற புதிய கணிதத்
துறையில் தொடர்ந்து ஆராய்ந்து வருகின்றனர்.

\end{document}
